\begin{abstract}
Creative generation often collapses toward safe and stereotyped responses, but the cause of that collapse may be architectural as well as decoding-driven. We study whether LayerNorm contributes to this effect by comparing vanilla \gpttwo against the released LN-free checkpoint \lnfreegpt under matched inference settings. Our evaluation combines 20 creative prompts from \noveltybench and 62 validation walls from \ocw, so we can measure both surface diversity and hidden-state association geometry. Removing LayerNorm increases semantic spread across repeated generations, with embedding dispersion rising from 0.803 to 0.831 ($\Delta=+0.0285$, FDR-adjusted $p=0.024$). It also produces a much broader internal geometry on \ocw: anisotropy drops from 0.992 to 0.964 ($\Delta=-0.0273$, FDR-adjusted $p=8.5\times 10^{-26}$) and effective rank rises from 1.35 to 2.56 ($\Delta=+1.21$, FDR-adjusted $p=1.5\times 10^{-11}$). However, these internal gains do not translate into reliable surface-level novelty. The LN-free model slightly reduces \texttt{distinct\_1} and \texttt{distinct\_2}, slightly increases lexical overlap, and does not significantly improve clue-group separability. The main conclusion is that LayerNorm compresses representation geometry in this GPT-2-family setting, but removing it alone is not enough to produce clearly more novel creative outputs.
\end{abstract}
